% ======================================================================
% Discussion.
% Leads with what the space changes in practice, then treats the three
% findings as evidence that it was worth building, then the construction
% principle, then the limits.
% ======================================================================
\section{Discussion}

What changes when the public record is held as one structurally resolved space is
first of all practical. A screening campaign that used to begin with account
requests, bulk downloads and format conversion begins instead with a query, and
the answer carries with it the source of every number in it. A crystal is one row
rather than six, so counting is possible; a property value names the calculation
behind it, so comparison is possible; and the licence position of every value is
a column rather than a matter to be re-established. This is infrastructure rather
than discovery, and its value is measured by what it removes from everyone else's
work.

It also makes a class of question answerable that was not. A repository can
report on its own contents; only a resolved union can report on the field. How
many distinct crystals are actually known, as opposed to how many records exist.
Which groups have already computed a given compound, and whether they agreed.
Whether a training set and a test set drawn from two providers share the same
crystals. None of these is a property of any member, and none can be recovered
afterwards from a corpus assembled by concatenation.

Two of the answers were not what we expected, and both bear on how the corpus
should be used. Independent sources reporting the same formation energy for the
same crystal agree within 50~meV\,atom$^{-1}$ only 60.5\% of the time, and the
least concordant tenth differ by margins several times the thresholds on which
stability screening turns. Because the discrepancy is a systematic offset rather
than scatter, it does not average away with sample size, and a corpus that merged
the contributing columns would hand a model the offset to learn as though it were
physics. The remedy is not to choose a preferred provider but to keep the
provenance, which costs one column per value and turns the disagreement into an
observable. The complementary caution is equally practical: a pair of sources
agreeing to the last decimal is usually carrying one calculation twice, so
agreement between repositories should be read as evidence of shared provenance
until shown otherwise. The second answer concerns what the corpus does not
contain. Coverage of the property space follows the price of the calculation, not
the value of the answer, and within the expensive tier it tracks whether a
dedicated high-throughput campaign happened to exist. The resulting landscape is a
record of where the community spent computer time, which is a different thing
from a record of what is worth knowing.

That landscape is also a well-posed learning problem, and the space supplies what
such a problem needs: 497 candidate tasks once each property name is split by
level of theory, 162 of them realised as training sets over 93 million labels,
deduplicated structures to split on, labels that retain the quantity they denote
and, on the entries with two independent sources, a measurement of the floor
below which a prediction cannot be distinguished from the disagreement of its own
training data. Section~\ref{sec:completion} also shows how easily that problem is
posed wrongly. Roughly two thirds of the tasks are defined on less than half the
corpus, and asking each of them about every crystal would generate 2.19 billion
numbers of which more than half would be undefined rather than merely uncertain.
Enforcing the domains is not a loss of coverage; it is the difference between a
completed property layer and a plausible-looking one.

The reason to complete it goes beyond adding numbers to a table. Materials
discovery runs on two routes, and the state of the public record has kept them
apart. Forward screening searches the known set for a material that already has
the property wanted; inverse design generates a structure to order. Forward
screening is habitually treated as bounded by how many materials are known, but
the corpus shows that it is bounded instead by how few of them have been
characterised: no property requiring a calculation of its own reaches a fiftieth
of the 13.5 million crystals, so a screen conducted over the public record ranks
the materials somebody chose to compute rather than the materials that exist, and
its apparent frontier is an artefact of where computer time was spent. Predicting
the missing entries across the whole corpus, inside the domains where each
property is defined, replaces that artefact with something closer to the real
frontier and answers a question the field has not been able to pose: how much of
the known record has been passed over for want of a calculation rather than for
want of merit. The two possible answers are both consequential. If the yield is
large, forward screening has been operating far inside its own reach and a
substantial number of already-synthesisable materials are waiting to be
recognised. If it is small, the reach of the known set is close to exhausted, and
further progress has to come from designing structures rather than finding them.
The same corpus serves that second route, because a generative model needs what
it holds: structures deduplicated across the whole published record, so that a
claim of novelty is tested against everything known rather than against one
repository, and property labels that name the quantity they denote, so that
generation can be conditioned on a defined target rather than on a column heading.
Standardising structure and property data at this scale is what makes the two
routes comparable, and eventually complementary, rather than separate literatures.

One property of the construction is worth separating from the results it
produced. Admission and matching are the same mechanism seen from two sides: a
record that matches nothing is not an error to be logged but the corpus's next
member, and once admitted it becomes part of the criterion by which the next
record is judged. Over the full assembly, 15.9\% of matches resolved onto
crystals that assembly itself had added, approaching unity for chemistries the
starting pool did not cover. The same relation reappears one level up in the
completion stage, where seventeen tasks can only be gated once earlier models
have supplied the property their domain is defined by. A space built this way has
no final state: each round of ingestion enlarges what the next round can
recognise, and each round of prediction enlarges the domain over which the next
round is defined. That is what makes it maintainable as its sources keep moving,
which a merged snapshot is not.

Several limitations bound these conclusions. Structural fingerprinting has a
domain of validity, and enumerated framework libraries fall outside it; the
composite key rescues the two families we encountered but is a remedy for a known
failure rather than a general solution. Agreement can be measured only where two
providers report the same quantity under comparable methods, which the method
strings make checkable but which holds for a minority of the occupied matrix, so
the disagreement reported here is measured on the part of the corpus most likely
to be well behaved. The applicability gate is structural rather than geometric: a
porous domain is enforced through family labels rather than a computed
pore-limiting diameter, which is the stricter test and remains to be run over the
whole corpus. Licence fragmentation persists downstream, since four sources
contributing 26.0\% of all values cannot be redistributed, so the public artifact
and the locally built one are not the same object and any published statistic has
to say which it describes. And no value here was recomputed. Every number is as
its source reported it, and the union inherits whatever errors those sources
contain; what it adds is the ability to see when two of them disagree.
